% Regression: frame-readable wire records reserve the faces their ink uses.
\documentclass{standalone}
\usepackage{tenkz}
\makeatletter
\ExplSyntaxOn
\file_input:n { tenkz-kernel.code.tex }
\int_new:N \l__tenkz_test_frame_side_count_int
\int_new:N \l__tenkz_test_frame_route_count_int
\cs_new_eq:NN
  \__tenkz_test_frame_side_auto:nN \__tenkz_kernel_r_label_auto_side:nN
\cs_new_eq:NN \__tenkz_test_frame_beadpath:nnn \tenkz@string@beadpath
\cs_set_protected:Npn \tenkz@string@beadpath #1#2#3
  {
    \str_if_eq:nnT {#3} {tenkz-label-route-point}
      {
        \int_gincr:N \l__tenkz_test_frame_route_count_int
        \str_case:nnF {#1}
          {
            {tenkz-bead-cup-1-2}   { }
            {tenkz-bead-wrap-1}     { }
          }
          { \errmessage{closure~occupancy~did~not~read~a~frozen~bead~route} }
      }
    \__tenkz_test_frame_beadpath:nnn {#1} {#2} {#3}
  }
\cs_set_protected:Npn \__tenkz_kernel_r_label_auto_side:nN #1#2
  {
    \__tenkz_test_frame_side_auto:nN {#1} #2
    \int_gincr:N \l__tenkz_test_frame_side_count_int
    \__tenkz_model_get:nnN {#1} {label} \l_tmpa_tl
    \str_case:VnF \l_tmpa_tl
      {
        {P}
          {
            \str_if_eq:VnF #2 {w}
              { \errmessage{far~port~slot~answered~the~wrong~face} }
          }
        {Q}
          {
            \str_if_eq:VnF #2 {w}
              { \errmessage{spanning~atom~name~answered~the~wrong~face} }
          }
        {R}
          {
            \str_if_eq:VnF #2 {n}
              { \errmessage{generated~closure~answered~the~wrong~face} }
          }
        {T}
          {
            \str_if_eq:VnF #2 {s}
              { \errmessage{closure~bead~answered~the~wrong~face} }
          }
        {U}
          {
            \str_if_eq:VnF #2 {n}
              { \errmessage{closure~corner~lost~one~of~its~two~faces} }
          }
      }
      { \errmessage{unexpected~label~in~frame-cell~fixture} }
  }
\ExplSyntaxOff
\makeatother
\begin{document}
\tenkzkernel
% The far port lies on the second row of a spanning atom.  The horizontal
% wire therefore takes the labelled atom's east face.
\begin{tenkz}[lattice={3x3}, bonds=none]
  \tn[at=(2,1), name=P, skin=dot,
      ports={90:virtual,270:virtual}]{P}
  \tn[at=(1,3), name=A, wires=2,
      ports={180@2:virtual}]{}
  \tn[at=(1,1), name=N]{} \tn[at=(3,1), name=S]{}
  \tnwire{P.90}{N} \tnwire{P.270}{S} \tnwire{P}{A.180@2}
\end{tenkz}
% The bare name denotes the middle of the three-row rectangle, not its
% anchor cell.  Its horizontal wire therefore takes the east face.
\begin{tenkz}[lattice={3x3}, bonds=none]
  \tn[at=(1,2), name=Q, skin=dot, wires=3,
      ports={90:virtual,270:virtual}]{Q}
  \tn[at=(1,1), name=QN]{}
  \tn[at=(3,1), name=QS]{}
  \tn[at=(2,3), name=QE]{}
  \tnwire{Q.90}{QN} \tnwire{Q.270}{QS} \tnwire{Q}{QE}
\end{tenkz}
% A one-sided south trace leaves the boundary atom through its south face.
\begin{tenkz}[lattice={1x1}, bonds=none, south=trace]
  \tn[at=(1,1), skin=dot]{R}
\end{tenkz}
% A bead a quarter of the way round an east cup sees the saved curved route,
% not the vertical chord between the cup's boundary sites.  The transverse
% probe puts a declared crossing at the bead and cuts the live cup route; the
% occupancy reading must retain the frozen, unsurgered route used to place T.
\begin{tenkz}[rows={wire,wire}, cols=1, bonds=none, east=cup]
  \tn[at=(1,1)]{} \\ \tn[at=(2,1)]{}
  \tn[at=on cup-1-2 0.25, name=T, skin=dot]{T}
  \tnwire[kind=string, name=probe, crossing=over,
    cross={over at crossing of probe and cup-1-2}]
    {1 n of T}{1 s of T}
\end{tenkz}
% At the first corner of a flat trace, the incoming and outgoing pieces take
% two adjacent faces.  Reading only their common secant would call the corner
% oblique and leave its south face falsely free.
\begin{tenkz}[lattice={1x1}, bonds=none, west=trace, east=trace]
  \tn[at=(1,1)]{}
  \tn[at=on wrap-1 0.19449155, skin=dot]{U}
\end{tenkz}
\ExplSyntaxOn
\int_compare:nNnF { \l__tenkz_test_frame_side_count_int } = {5}
  { \errmessage{frame-cell~fixture~did~not~station~five~labels} }
\int_compare:nNnF { \l__tenkz_test_frame_route_count_int } = {6}
  { \errmessage{closure~occupancy~did~not~sample~two~frozen~routes} }
\ExplSyntaxOff
\end{document}
